\chapter{ETL Multi-Sumber dengan Python}
\label{modul:etl-multi-sumber}

\section{Identitas Modul}

\begin{identitasmodul}
\textbf{Sumber materi}    & Pertemuan 9 \\
\textbf{CPMK}             & CPMK-092 \\
\textbf{Minggu pelaksanaan} & Minggu ke-10 \\
\textbf{Durasi tatap muka} & 120 menit \\
\textbf{Estimasi tugas}   & 5--6 jam di luar sesi \\
\textbf{Moda kerja}       & Individual \\
\textbf{Perangkat}        & PostgreSQL 17, ekstensi \perintah{postgres\_fdw}, dan Python 3.11+ \\
\textbf{Dataset}          & Tiga sumber operasional NusaMart yang sengaja tidak sepakat \\
\textbf{Skrip pendukung}  & \berkas{etl/src/} beserta \berkas{sql/modul-07/check.sql} \\
\textbf{Luaran}           & \berkas{pipeline.py}, \berkas{log-muat.md}, dan laporan konflik skema \\
\end{identitasmodul}

\slideref{9}

\section{Capaian dan Indikator Keberhasilan}

Mahasiswa mampu mengekstrak beberapa sumber, mengenali konflik, menstandardisasi
data, melakukan object identification, serta memuat SCD-2 dan fakta secara
idempoten. Baris tanpa pasangan dimensi harus diarahkan ke baris unknown.

\begin{capaian}
Pada akhir modul ini mahasiswa mampu:
\begin{enumerate}[leftmargin=1.4em]
  \item menghubungkan basis data sumber lain secara langsung dengan
        \perintah{postgres\_fdw}, tanpa mengekspor berkas lebih dahulu;
  \item mengenali dan menggolongkan konflik antarsumber --- penamaan, satuan,
        pengodean, representasi, dan identitas;
  \item menulis pipeline Python yang memisahkan tahap ekstraksi,
        standardisasi, dan pemuatan secara tegas;
  \item melakukan \istilah{object identification} sehingga satu pelanggan yang
        tercatat berbeda di dua sumber menjadi satu baris dimensi;
  \item menangani baris yang datang terlambat tanpa merusak histori dimensi;
  \item membuktikan pipeline bersifat \istilah{idempoten} --- dijalankan dua
        kali menghasilkan keadaan yang sama.
\end{enumerate}
\end{capaian}

Sampai Modul 6, seluruh data berasal dari satu sumber yang sudah rapi dan
sepakat dengan dirinya sendiri. Mulai modul ini persoalannya berbeda jenis:
tiga sistem yang dibangun tim berbeda pada waktu berbeda, yang menyebut hal
yang sama dengan nama berbeda --- dan kadang menyebut hal yang berbeda dengan
nama yang sama.

\section{Prasyarat dan Persiapan}

\subsection{Prasyarat}

\begin{itemize}
  \item Python dasar: fungsi, \perintah{dict}, \perintah{list comprehension},
        dan penanganan galat dengan \perintah{try};
  \item \perintah{pandas} tingkat dasar: membaca \perintah{DataFrame},
        menyaring, dan menggabungkan;
  \item koneksi basis data dari Python dengan \perintah{psycopg2} atau
        \perintah{SQLAlchemy};
  \item pemuatan SCD-2 dari Modul 3, termasuk aturan menutup versi lama sebelum
        membuka versi baru;
  \item konsep baris unknown dari Modul 5.
\end{itemize}

\subsection{Pre-lab}

\begin{enumerate}[leftmargin=1.4em]
  \item Aktifkan ekstensi pada basis data gudang:
\begin{lstlisting}[style=sqlgaya]
CREATE EXTENSION IF NOT EXISTS postgres_fdw;
\end{lstlisting}
  \item Pastikan ketiga basis data sumber dapat dijangkau. Ketiganya berada
        pada container \perintah{postgres\_source} yang sama, sebagai tiga
        basis data terpisah.
\begin{lstlisting}[style=shellgaya]
docker compose exec -T postgres_source \
  psql -U praktikum -l | grep nusamart
\end{lstlisting}
  \item Siapkan lingkungan Python dan salin berkas konfigurasi:
\begin{lstlisting}[style=shellgaya]
python3 -m venv .venv && source .venv/bin/activate
python -m pip install -r requirements.txt
cp etl/config/.env.example .env
\end{lstlisting}
  \item Siapkan jawaban berikut pada \berkas{modul-07-etl/pre-lab.md}:
  \begin{enumerate}[label=\alph*.,leftmargin=1.4em]
    \item Sumber kedua menyimpan berat produk dalam kilogram, sumber pertama
          dalam gram. Bila keduanya dimuat apa adanya, kesalahan apa yang
          muncul, dan mengapa ia tidak akan menimbulkan galat?
    \item Pelanggan yang sama tercatat sebagai ``Siti Nurhaliza'' di satu
          sumber dan ``SITI NUR HALIZA'' di sumber lain. Sebutkan dua cara
          memutuskan bahwa keduanya orang yang sama, beserta risiko
          masing-masing.
    \item Apa arti ``pipeline yang idempoten'', dan bagaimana Anda akan
          mengujinya?
  \end{enumerate}
\end{enumerate}

\section{Konsep Dasar}

\subsection{Extract, Transform, dan Load}

Ketiga tahap memiliki batas tanggung jawab yang tegas, dan kedisiplinan menjaga
batas itulah yang membuat pipeline dapat diperbaiki ketika rusak.

\begin{center}
\small
\begin{tabular}{@{}p{0.16\textwidth}>{\raggedright\arraybackslash}p{0.36\textwidth}>{\raggedright\arraybackslash}p{0.40\textwidth}@{}}
\toprule
\textbf{Tahap} & \textbf{Tugasnya} & \textbf{Yang bukan tugasnya} \\
\midrule
Extract & Menyalin sumber apa adanya ke \perintah{staging} &
  Membersihkan, menyaring, atau memperbaiki \\
Transform & Menstandardisasi satuan, kode, dan identitas &
  Memutuskan versi dimensi \\
Load & Memuat dimensi SCD-2 dan fakta &
  Mengubah nilai measure \\
\bottomrule
\end{tabular}
\end{center}

\begin{peringatan}
Godaan terbesar adalah membersihkan data pada tahap ekstraksi --- menyaring
baris rusak, membetulkan satuan, membuang duplikat --- karena tampak lebih
efisien. Jangan. Begitu data rusak tidak pernah sampai ke \perintah{staging},
Anda kehilangan kemampuan menghitung berapa banyak yang rusak, dan
rekonsiliasi Modul 10 tidak lagi mungkin. \textbf{Staging menampung, bukan
menolak.}
\end{peringatan}

\subsection{Foreign Data Wrapper dan Staging}

\perintah{postgres\_fdw} membuat tabel pada basis data lain tampak seperti
tabel biasa. Tidak ada ekspor CSV, tidak ada berkas perantara, dan tidak ada
langkah manual yang bisa terlewat.

\begin{konsep}{Empat langkah menyiapkan FDW}
\begin{enumerate}[leftmargin=1.4em]
  \item \perintah{CREATE SERVER} --- mendaftarkan alamat basis data seberang.
  \item \perintah{CREATE USER MAPPING} --- menyatakan kredensial yang dipakai.
  \item \perintah{IMPORT FOREIGN SCHEMA} --- membuat definisi tabel asing.
  \item \perintah{SELECT} --- dipakai seperti tabel biasa.
\end{enumerate}
\end{konsep}

Kemudahan itu ada batasnya. Setiap query atas tabel asing menempuh jaringan,
dan tidak seluruh operasi dapat didorong ke sisi seberang. Karena itu FDW
dipakai untuk \emph{memindahkan sekali} ke \perintah{staging}, bukan sebagai
sumber yang di-query berulang kali oleh laporan.

\subsection{Standardisasi dan Object Identification}

Konflik antarsumber muncul dalam beberapa jenis, dan mengenalinya lebih dahulu
mempercepat penyelesaiannya.

\begin{center}
\small
\begin{tabular}{@{}p{0.20\textwidth}>{\raggedright\arraybackslash}p{0.34\textwidth}>{\raggedright\arraybackslash}p{0.38\textwidth}@{}}
\toprule
\textbf{Jenis konflik} & \textbf{Wujudnya} & \textbf{Contoh pada NusaMart} \\
\midrule
Penamaan & Nama kolom berbeda untuk hal sama &
  \perintah{produk\_id} lawan \perintah{item\_code} \\
Satuan & Besaran sama, satuan berbeda &
  Berat dalam gram lawan kilogram \\
Pengodean & Nilai sama, kode berbeda &
  Kategori \perintah{MNM} lawan \perintah{BEV} \\
Representasi & Format berbeda &
  Tanggal, huruf besar-kecil, tanda baca \\
Identitas & Entitas sama, catatan berbeda &
  Satu pelanggan pada dua sumber \\
\bottomrule
\end{tabular}
\end{center}

Empat jenis pertama diselesaikan dengan aturan yang dapat ditulis sekali.
Jenis kelima berbeda: ia memerlukan \emph{keputusan}, dan keputusan itu bisa
salah.

\begin{konsep}{Object identification dalam empat langkah}
\begin{enumerate}[leftmargin=1.4em]
  \item \textbf{Normalisasi.} Huruf kecil semua, buang tanda baca, rapatkan
        spasi ganda, buang gelar dan sapaan.
  \item \textbf{Blocking.} Bandingkan hanya calon yang masuk akal --- misalnya
        yang kotanya sama. Tanpa ini, 100 ribu pelanggan menuntut lima miliar
        perbandingan.
  \item \textbf{Penilaian.} Hitung kemiripan nama, lalu gabungkan dengan
        kecocokan atribut lain seperti tanggal lahir.
  \item \textbf{Ambang dan keputusan.} Di atas ambang: satu orang. Di bawah
        ambang: dua orang. Di sekitar ambang: \emph{ditinjau manusia}.
\end{enumerate}
\end{konsep}

\begin{peringatan}
Ambang yang terlalu longgar menggabungkan dua orang berbeda menjadi satu ---
seluruh transaksi mereka bercampur, dan kesalahan ini nyaris mustahil
ditemukan kembali. Ambang yang terlalu ketat membiarkan satu orang tercatat
dua kali --- kesalahan yang lebih mudah ditemukan dan diperbaiki.

Bila ragu, pilih yang terlalu ketat. Menggabungkan dua baris nanti jauh lebih
mudah daripada memisahkan transaksi yang sudah tercampur.
\end{peringatan}

Sesudah dua catatan dinyatakan sebagai satu orang, muncul pertanyaan
berikutnya: nilai mana yang dipakai? Inilah \istilah{survivorship}. Aturannya
harus ditulis --- misalnya ``sumber pertama menang untuk alamat, catatan
terbaru menang untuk nomor telepon'' --- bukan diputuskan diam-diam oleh
urutan pemrosesan.

\subsection{Late-arriving Data dan Unknown Member}

Sebagian transaksi dari sumber ketiga baru sampai tiga hari sesudah kejadian.
Ini menimbulkan dua persoalan yang berbeda dan sering tertukar.

\begin{center}
\small
\begin{tabular}{@{}p{0.24\textwidth}>{\raggedright\arraybackslash}p{0.32\textwidth}>{\raggedright\arraybackslash}p{0.36\textwidth}@{}}
\toprule
\textbf{Persoalan} & \textbf{Wujudnya} & \textbf{Penanganan} \\
\midrule
\istilah{Late-arriving fact} & Transaksi lama baru sampai hari ini &
  Cari versi dimensi yang berlaku pada tanggal \emph{transaksi}, bukan hari
  ini \\
\istilah{Late-arriving dimension} & Fakta menyebut produk yang belum ada di
  dimensi & Arahkan ke baris unknown, atau buat versi awal beruntung \\
\bottomrule
\end{tabular}
\end{center}

Yang pertama sudah Anda tangani sejak Modul 3: pencarian kunci berbasis rentang
tanggal, bukan \perintah{baris\_kini}. Justru karena itulah bentuk join tersebut
dipilih sejak awal --- dan pada modul ini manfaatnya baru benar-benar terasa.

\subsection{Idempotensi dan Audit Pemuatan}

\begin{konsep}{Pipeline yang idempoten}
Menjalankan pipeline dua kali atas masukan yang sama menghasilkan keadaan basis
data yang \emph{sama persis} dengan menjalankannya sekali. Tidak ada baris
kembar, tidak ada versi dimensi tambahan, dan tidak ada measure yang terhitung
dua kali.
\end{konsep}

Ini bukan kemewahan. Pipeline gagal di tengah jalan adalah hal biasa, dan
satu-satunya penanganan yang aman adalah menjalankannya ulang. Pipeline yang
tidak idempoten membuat setiap kegagalan menuntut pembersihan manual.

Tiga teknik yang dipakai modul ini:

\begin{enumerate}[leftmargin=1.4em]
  \item \textbf{Kunci alami pada staging} --- pemuatan ulang menimpa, bukan
        menambah, dengan \perintah{ON CONFLICT DO UPDATE}.
  \item \textbf{Hapus-lalu-muat per batch} --- fakta satu tanggal dihapus lebih
        dahulu sebelum dimuat ulang.
  \item \textbf{Perbandingan atribut sebelum membuat versi} --- SCD-2 hanya
        membuat versi baru bila atributnya benar-benar berubah, sebagaimana
        Modul 3.
\end{enumerate}

Setiap jalannya pipeline mencatat dirinya pada tabel audit: berapa baris
dibaca, dimuat, dan ditolak per sumber. Tabel ini menjadi bahan
\istilah{audit dimension} pada Modul 8 dan rekonsiliasi pada Modul 10.

\section{Studi Kasus dan Protokol Praktikum}

Dua sumber tambahan sengaja memakai gram dan kilogram, kode kategori berbeda,
ejaan pelanggan tidak konsisten, serta transaksi terlambat tiga hari.

\begin{center}
\small
\begin{tabular}{@{}p{0.24\textwidth}p{0.16\textwidth}>{\raggedright\arraybackslash}p{0.52\textwidth}@{}}
\toprule
\textbf{Basis data} & \textbf{Peran} & \textbf{Ciri yang menyulitkan} \\
\midrule
\perintah{nusamart\_oltp} & Sumber utama & Acuan penamaan; berat dalam gram \\
\perintah{nusamart\_barat} & Hasil akuisisi & Nama kolom berbahasa Inggris;
  berat dalam kilogram; kode kategori sendiri \\
\perintah{nusamart\_timur} & Wilayah timur & Ejaan pelanggan tidak konsisten;
  transaksi sampai terlambat tiga hari \\
\bottomrule
\end{tabular}
\end{center}

Dimensi pelanggan yang sengaja ditunda sejak Modul 2 akhirnya dibangun pada
modul ini. Penundaannya beralasan: membangun \perintah{dim\_pelanggan} dari
satu sumber tidak mengajarkan apa pun, sedangkan membangunnya dari tiga sumber
yang tidak sepakat adalah inti persoalan integrasi data.

\begin{catatanasisten}
Ketiga basis data berada pada container \perintah{postgres\_source} yang sama,
sehingga \berkas{docker-compose.yml} tidak bertambah layanannya. Yang berubah
hanya skrip seed: ia kini membangkitkan tiga basis data dengan
\perintah{--seed} yang sama tetapi aturan penamaan dan satuan yang berbeda.

Pastikan konflik yang ditanam \emph{terdokumentasi} pada berkas asisten, tetapi
\emph{tidak} dibagikan ke mahasiswa. Menemukan konflik adalah pekerjaan
Bagian~B.
\end{catatanasisten}

\section{Alur Praktikum 120 Menit}

\begin{alurwaktu}
0--15 & Demonstrasi FDW dan satu konflik skema. \\
15--95 & Kerja terbimbing: profiling, transformasi, identity resolution, dan load. \\
95--110 & Checkpoint hasil rekonsiliasi dan penanganan unknown member. \\
110--120 & Penjelasan pembuktian idempotensi. \\
\end{alurwaktu}

\section{Prosedur Praktikum Terbimbing}

\subsection{Bagian A: Menghubungkan Sumber dengan FDW}

\langkah{A.1 Mendaftarkan server dan pemetaan pengguna}{10 menit}

\begin{lstlisting}[style=sqlgaya]
CREATE SERVER src_barat FOREIGN DATA WRAPPER postgres_fdw
  OPTIONS (host 'postgres_source', port '5432', dbname 'nusamart_barat');

CREATE USER MAPPING FOR CURRENT_USER SERVER src_barat
  OPTIONS (user 'praktikum', password 'ganti-nilai-ini');

CREATE SCHEMA IF NOT EXISTS src_barat;
IMPORT FOREIGN SCHEMA public FROM SERVER src_barat INTO src_barat;
\end{lstlisting}

\begin{catatan}
Perhatikan \perintah{host 'postgres\_source'} dan \perintah{port '5432'} ---
nama layanan dan port \emph{di dalam} jaringan Docker, bukan
\perintah{localhost:5433} yang Anda pakai dari laptop. Kekeliruan ini
menimbulkan pesan \emph{could not connect to server} yang membingungkan karena
koneksi dari DBeaver jelas berhasil.
\end{catatan}

Ulangi untuk \perintah{src\_timur}, lalu uji akses:

\begin{lstlisting}[style=sqlgaya]
SELECT COUNT(*) FROM src_barat.item;
SELECT COUNT(*) FROM src_timur.pelanggan;
\end{lstlisting}

\subsection{Bagian B: Memprofilkan Konflik Sumber}

\langkah{B.1 Membandingkan struktur ketiga sumber}{10 menit}

\begin{lstlisting}[style=sqlgaya]
SELECT table_schema, table_name, column_name, data_type
FROM   information_schema.columns
WHERE  table_schema IN ('src_barat', 'src_timur')
ORDER  BY table_schema, table_name, ordinal_position;
\end{lstlisting}

Susun tabel pemetaan kolom pada \berkas{laporan-konflik.md}: untuk setiap
konsep bisnis, nama kolomnya pada masing-masing sumber.

\langkah{B.2 Menemukan konflik satuan dan pengodean}{10 menit}

Konflik satuan tidak terlihat dari nama kolom. Ia terlihat dari
\emph{sebaran nilai}.

\begin{lstlisting}[style=sqlgaya]
-- Bandingkan besaran yang seharusnya setara.
SELECT 'oltp (gram)' AS sumber,
       MIN(berat_gram), ROUND(AVG(berat_gram), 2), MAX(berat_gram)
FROM   src_oltp.produk
UNION ALL
SELECT 'barat (?)', MIN(weight_kg), ROUND(AVG(weight_kg), 2), MAX(weight_kg)
FROM   src_barat.item;

-- Kode kategori yang dipakai masing-masing sumber.
SELECT 'oltp' AS sumber, kode_kategori, COUNT(*)
FROM   src_oltp.kategori GROUP BY 1, 2
UNION ALL
SELECT 'barat', category_code, COUNT(*)
FROM   src_barat.item_category GROUP BY 1, 2
ORDER  BY 1, 2;
\end{lstlisting}

Rata-rata yang berselisih sekitar seribu kali adalah tanda konflik satuan.
Catat temuan ini beserta angkanya --- inilah jenis bukti yang diminta.

\langkah{B.3 Menemukan konflik identitas}{10 menit}

\begin{lstlisting}[style=sqlgaya]
-- Calon pelanggan yang sama, terlihat dari nama yang dinormalkan.
SELECT lower(regexp_replace(nama, '[^a-zA-Z ]', '', 'g')) AS nama_normal,
       COUNT(*) AS jumlah_catatan,
       string_agg(DISTINCT sumber, ', ') AS muncul_di
FROM ( SELECT nama, 'oltp'  AS sumber FROM src_oltp.pelanggan
       UNION ALL
       SELECT cust_name, 'barat' FROM src_barat.customer
       UNION ALL
       SELECT nama, 'timur'  FROM src_timur.pelanggan ) g
GROUP  BY 1
HAVING COUNT(*) > 1
ORDER  BY jumlah_catatan DESC
LIMIT  20;
\end{lstlisting}

Lengkapi \berkas{laporan-konflik.md} dengan sekurang-kurangnya \textbf{lima}
konflik, masing-masing disertai jenisnya, buktinya berupa angka, dan aturan
penyelesaian yang Anda usulkan.

\begin{luaran}
\berkas{modul-07-etl/laporan-konflik.md} --- tabel pemetaan kolom, lima konflik
beserta bukti angka, dan aturan penyelesaian untuk masing-masing.
\end{luaran}

\subsection{Bagian C: Menulis Transformasi Python}

\langkah{C.1 Menyusun kerangka pipeline}{10 menit}

Kerangka tersedia pada \berkas{etl/src/}. Perhatikan pemisahan tahapnya ---
inilah yang membuat pipeline dapat diuji sebagian.

\begin{lstlisting}[style=pythongaya]
# pipeline.py
from extract   import ekstrak_semua_sumber
from transform import standardisasi, identifikasi_pelanggan
from load      import muat_dimensi, muat_fakta
from audit     import mulai_batch, selesai_batch

def jalankan(tanggal_batch):
    batch = mulai_batch(tanggal_batch)
    mentah = ekstrak_semua_sumber(batch)        # -> staging, apa adanya
    bersih = standardisasi(mentah)              # satuan, kode, format
    bersih = identifikasi_pelanggan(bersih)     # -> pelanggan_master_id
    muat_dimensi(bersih, batch)                 # SCD-2
    muat_fakta(bersih, batch)                   # + baris unknown
    selesai_batch(batch)
\end{lstlisting}

\langkah{C.2 Menstandardisasi satuan dan kode}{15 menit}

\begin{lstlisting}[style=pythongaya]
# Aturan konversi ditulis sebagai data, bukan sebagai percabangan if.
# Menambah sumber keempat berarti menambah satu baris, bukan mengubah kode.
KONVERSI_BERAT = {
    "oltp":  1.0,      # sudah gram
    "barat": 1000.0,   # kilogram -> gram
    "timur": 1.0,
}

PETA_KATEGORI = {
    ("barat", "BEV"): "MNM",
    ("barat", "SNK"): "MKN",
    ("timur", "DRK"): "MNM",
}

def standardisasi(df):
    df["berat_gram"] = df.apply(
        lambda r: r["berat"] * KONVERSI_BERAT[r["sumber"]], axis=1)

    df["kode_kategori"] = df.apply(
        lambda r: PETA_KATEGORI.get((r["sumber"], r["kode_kategori"]),
                                    r["kode_kategori"]), axis=1)
    return df
\end{lstlisting}

\begin{catatan}
Menulis aturan sebagai \perintah{dict} dan bukan sebagai rangkaian
\perintah{if} bukan sekadar selera. Aturan yang berupa data dapat dicetak,
diperiksa, diuji, dan --- pada Modul 8 --- didokumentasikan sebagai metadata.
Aturan yang tersembunyi di dalam percabangan hanya dapat dibaca oleh yang
menulisnya.
\end{catatan}

\langkah{C.3 Mengidentifikasi pelanggan lintas sumber}{15 menit}

\begin{lstlisting}[style=pythongaya]
import re
from difflib import SequenceMatcher

GELAR = {"bapak", "ibu", "bpk", "sdr", "sdri", "h", "hj", "drs", "ir"}
AMBANG = 0.88

def normalkan(nama):
    n = re.sub(r"[^a-z ]", " ", nama.lower())
    kata = [k for k in n.split() if k not in GELAR]
    return " ".join(kata)

def mirip(a, b):
    return SequenceMatcher(None, a, b).ratio()

def identifikasi_pelanggan(catatan):
    """Blocking menurut kota, lalu penilaian kemiripan nama."""
    master, ragu = {}, []
    for kota, kelompok in kelompokkan(catatan, kunci="kota"):
        wakil = []                       # (nama_normal, master_id)
        for c in kelompok:
            n = normalkan(c["nama"])
            cocok = [(mirip(n, wn), mid) for wn, mid in wakil]
            skor, mid = max(cocok, default=(0.0, None))

            if skor >= AMBANG:
                master[c["id_sumber"]] = mid
            elif skor >= AMBANG - 0.08:
                ragu.append((c, skor, mid))      # ditinjau manusia
            else:
                mid = master_id_baru()
                wakil.append((n, mid))
                master[c["id_sumber"]] = mid
    return master, ragu
\end{lstlisting}

Simpan daftar \perintah{ragu} ke \berkas{tinjau-manual.csv}. Jumlah barisnya
adalah bagian dari laporan --- pipeline yang tidak pernah ragu hampir pasti
ambangnya terlalu longgar.

\begin{peringatan}
Blocking menurut kota mempercepat proses, tetapi juga berarti pelanggan yang
sama dengan kota berbeda \emph{tidak akan pernah} dibandingkan. Ini pertukaran
yang disengaja dan harus Anda catat sebagai keterbatasan pada laporan, bukan
disembunyikan.
\end{peringatan}

\subsection{Bagian D: Memuat Dimensi dan Fakta}

\langkah{D.1 Membangun dan memuat dim\_pelanggan}{10 menit}

\begin{lstlisting}[style=sqlgaya]
CREATE TABLE gudang.dim_pelanggan (
  pelanggan_key   INTEGER GENERATED ALWAYS AS IDENTITY PRIMARY KEY,
  pelanggan_id    INTEGER      NOT NULL,   -- master id hasil identifikasi
  nama            VARCHAR(120) NOT NULL,
  kota            VARCHAR(60),
  segmen          VARCHAR(30),
  sumber_asal     VARCHAR(10)  NOT NULL,   -- sumber yang memenangkan nilai
  mulai_berlaku   DATE    NOT NULL,
  selesai_berlaku DATE    NOT NULL,
  baris_kini      BOOLEAN NOT NULL,
  versi           SMALLINT NOT NULL
);

CREATE UNIQUE INDEX uq_dim_pelanggan_kini
  ON gudang.dim_pelanggan (pelanggan_id) WHERE baris_kini;

ALTER TABLE gudang.dim_pelanggan ADD CONSTRAINT ex_dim_pelanggan_rentang
  EXCLUDE USING gist (
    pelanggan_id WITH =,
    daterange(mulai_berlaku, selesai_berlaku, '[)') WITH &&);
\end{lstlisting}

Struktur SCD-2-nya sama persis dengan \perintah{dim\_produk} Modul 3, termasuk
kedua penjaganya. Jangan lupa baris unknown berkunci $-1$ sebagaimana Modul 5.

\langkah{D.2 Memuat fakta dengan penanganan keterlambatan}{15 menit}

\begin{lstlisting}[style=pythongaya]
SQL_MUAT_FAKTA = """
INSERT INTO gudang.fakta_penjualan (
  tanggal_key, produk_key, toko_key, pelanggan_key, flag_key,
  nomor_struk, transaksi_id, kuantitas, harga_satuan, diskon, subtotal)
SELECT dt.tanggal_key,
       COALESCE(dp.produk_key, -1),        -- baris unknown, bukan NULL
       COALESCE(dtk.toko_key, -1),
       COALESCE(dpl.pelanggan_key, -1),
       COALESCE(fl.flag_key, -1),
       s.nomor_struk, s.transaksi_id,
       s.kuantitas, s.harga_satuan, s.diskon, s.subtotal
FROM       staging.penjualan_gabungan s
JOIN       gudang.dim_tanggal dt ON dt.tanggal = s.waktu_transaksi::DATE
-- LEFT JOIN, bukan JOIN: baris tanpa pasangan harus tetap masuk
LEFT JOIN  gudang.dim_produk dp
       ON  dp.produk_id = s.produk_id
       AND s.waktu_transaksi::DATE >= dp.mulai_berlaku
       AND s.waktu_transaksi::DATE <  dp.selesai_berlaku
LEFT JOIN  gudang.dim_pelanggan dpl
       ON  dpl.pelanggan_id = s.pelanggan_master_id
       AND s.waktu_transaksi::DATE >= dpl.mulai_berlaku
       AND s.waktu_transaksi::DATE <  dpl.selesai_berlaku
LEFT JOIN  gudang.dim_toko dtk ON dtk.toko_id = s.toko_id
LEFT JOIN  gudang.dim_transaksi_flag fl
       ON  fl.metode_bayar = s.metode_bayar AND fl.status = s.status
WHERE  s.batch_id = %(batch)s;
"""
\end{lstlisting}

Dua hal menentukan kebenaran perintah ini. \perintah{LEFT JOIN} memastikan
baris tanpa pasangan tetap masuk, dan \perintah{COALESCE(..., -1)} mengarahkan
kunci yang kosong ke baris unknown --- bukan membiarkannya \perintah{NULL}.
Pencarian dimensi tetap memakai rentang tanggal transaksi, sehingga transaksi
yang terlambat tiga hari tetap menunjuk versi yang berlaku saat ia terjadi.

\langkah{D.3 Mencatat audit pemuatan}{5 menit}

\begin{lstlisting}[style=sqlgaya]
CREATE TABLE gudang.audit_muat (
  batch_id       BIGINT      NOT NULL,
  tahap          VARCHAR(20) NOT NULL,
  sumber         VARCHAR(10) NOT NULL,
  baris_dibaca   BIGINT      NOT NULL,
  baris_dimuat   BIGINT      NOT NULL,
  baris_unknown  BIGINT      NOT NULL DEFAULT 0,
  mulai          TIMESTAMP   NOT NULL,
  selesai        TIMESTAMP,
  PRIMARY KEY (batch_id, tahap, sumber)
);
\end{lstlisting}

Setiap jalannya pipeline mengisi tabel ini. Selisih antara
\perintah{baris\_dibaca} dan \perintah{baris\_dimuat} wajib dapat dijelaskan.

\begin{luaran}
\berkas{modul-07-etl/pipeline.py} beserta modul pendukungnya,
\berkas{log-muat.md} berisi isi tabel audit untuk dua kali jalan, dan
\berkas{tinjau-manual.csv}.
\end{luaran}

\begin{checkpoint}
Asisten menjalankan \berkas{sql/modul-07/check.sql} dan memeriksa butir berikut:
\begin{ceklis}
  \item ketiga sumber terhubung lewat \perintah{postgres\_fdw};
  \item berat produk sudah dalam satuan yang sama, terbukti dari sebaran
        nilainya;
  \item kode kategori sudah terstandardisasi;
  \item pelanggan yang sama dari dua sumber menjadi \textbf{satu} baris
        dimensi;
  \item tidak ada foreign key fakta yang \perintah{NULL} --- seluruhnya
        menunjuk dimensi atau baris unknown;
  \item transaksi yang datang terlambat menunjuk versi dimensi yang berlaku
        saat transaksi terjadi, bukan versi kini;
  \item tabel \perintah{audit\_muat} terisi dan selisihnya dapat dijelaskan;
  \item \berkas{laporan-konflik.md} memuat sekurang-kurangnya lima konflik
        beserta bukti angka.
\end{ceklis}
\end{checkpoint}

\section{Latihan Mandiri di Kelas}

\latihan{Menambahkan satu konflik baru}

Sumber \perintah{nusamart\_timur} mencatat tanggal transaksi dalam zona waktu
setempat tanpa penanda zona, sedangkan sumber utama memakai UTC.

\begin{enumerate}[leftmargin=1.4em]
  \item Tulis query yang membuktikan adanya perbedaan ini dari data, bukan dari
        dugaan.
  \item Berapa banyak transaksi yang akan jatuh pada \emph{tanggal yang salah}
        bila perbedaan itu diabaikan?
  \item Tulis aturan penyelesaiannya, lalu jelaskan di tahap mana aturan itu
        seharusnya diterapkan --- extract, transform, atau load.
\end{enumerate}

\latihan{Menguji ambang identifikasi}

Jalankan ulang identifikasi pelanggan dengan tiga ambang: 0,80, 0,88, dan 0,95.

\begin{enumerate}[leftmargin=1.4em]
  \item Berapa jumlah pelanggan master yang dihasilkan masing-masing?
  \item Ambil satu pasangan yang tergabung pada 0,80 tetapi terpisah pada 0,95.
        Menurut Anda, mana yang benar?
  \item Ambang mana yang Anda pilih, dan bukti apa yang akan Anda minta sebelum
        menaikkan atau menurunkannya?
\end{enumerate}

\section{Tugas Individu}

\subsection{Pekerjaan Wajib}
Jalankan pipeline dua kali berturut-turut dan buktikan jumlah baris stabil serta
tidak terbentuk versi dimensi kembar.

Kumpulkan \berkas{bukti-idempotensi.md} yang memuat:

\begin{enumerate}[leftmargin=1.4em]
  \item jumlah baris setiap tabel gudang sesudah jalan pertama dan sesudah
        jalan kedua, disandingkan dalam satu tabel;
  \item jumlah versi dimensi per natural key sebelum dan sesudah jalan kedua,
        membuktikan tidak ada versi kembar terbentuk;
  \item isi \perintah{audit\_muat} untuk kedua jalan;
  \item bila ada selisih, penjelasan penyebabnya dan perbaikan yang Anda
        lakukan --- selisih yang dijelaskan lebih bernilai daripada selisih
        yang disembunyikan;
  \item satu paragraf tentang bagian pipeline mana yang paling rapuh terhadap
        pengulangan, beserta alasannya.
\end{enumerate}

\subsection{Pertanyaan Analisis}

\begin{enumerate}[leftmargin=1.4em]
  \item Ambang identifikasi Anda menggabungkan dua orang yang sebenarnya
        berbeda. Jelaskan bagaimana kesalahan ini akan tampak pada laporan
        penjualan, dan mengapa ia jauh lebih sulit ditemukan daripada
        kesalahan sebaliknya.
  \item Sumber kedua mengirim ulang seluruh data tiga bulan terakhir setiap
        malam, bukan hanya data yang berubah. Apa akibatnya bagi pipeline Anda,
        dan bagaimana Anda menanganinya tanpa membuat versi dimensi kembar?
  \item Pipeline Anda mengarahkan produk yang tidak dikenal ke baris unknown.
        Sebulan kemudian produk itu muncul di dimensi. Apakah baris fakta lama
        diperbaiki, dibiarkan, atau dimuat ulang? Uraikan konsekuensi ketiga
        pilihan.
\end{enumerate}

\section{Format Pengumpulan dan Checklist}

Kumpulkan pipeline, konfigurasi contoh tanpa kredensial, log muat, laporan
konflik, dan bukti dua kali eksekusi.

Seluruh berkas diletakkan pada \berkas{modul-07-etl/}.

\begin{ceklis}
  \item \berkas{pre-lab.md}
  \item \berkas{ddl-fdw.sql} dan \berkas{ddl-dim-pelanggan.sql}
  \item \berkas{laporan-konflik.md} --- lima konflik beserta bukti angka
  \item \berkas{pipeline.py} beserta modul pendukungnya
  \item \berkas{.env.example} --- \textbf{tanpa} kredensial sebenarnya
  \item \berkas{log-muat.md} dan \berkas{tinjau-manual.csv}
  \item \berkas{bukti-idempotensi.md}
  \item \berkas{jawaban-analisis.md}
\end{ceklis}

\begin{peringatan}
Berkas \berkas{.env} berisi kata sandi dan \textbf{tidak boleh} masuk
repositori. Periksa \berkas{.gitignore} sebelum melakukan commit. Kredensial
yang pernah ter-commit tetap ada pada riwayat Git meskipun berkasnya kemudian
dihapus.
\end{peringatan}

\section{Troubleshooting}

\begin{table}[htbp]
\centering
\small
\caption{Masalah yang lazim muncul pada Modul 7 dan penanganannya}
\label{tab:m7-troubleshooting}
\begin{tabular}{@{}p{0.30\textwidth}>{\raggedright\arraybackslash}p{0.62\textwidth}@{}}
\toprule
\textbf{Gejala} & \textbf{Penanganan} \\
\midrule
\perintah{could not connect to server} pada FDW &
  Alamat memakai \perintah{localhost:5433} milik laptop. Dari dalam jaringan
  Docker, alamatnya \perintah{postgres\_source:5432}. \\
\perintah{password authentication failed} pada FDW &
  Kata sandi pada \perintah{USER MAPPING} berbeda dengan
  \berkas{.env}. Perbarui pemetaannya dengan
  \perintah{ALTER USER MAPPING}. \\
\perintah{relation "src\_barat.item" does not exist} &
  \perintah{IMPORT FOREIGN SCHEMA} belum dijalankan, atau schema sumbernya
  bukan \perintah{public}. \\
Angka berat berselisih sekitar seribu kali &
  Konflik satuan gram lawan kilogram. Bukan kerusakan data --- terapkan
  konversi pada tahap transform. \\
Jumlah pelanggan master jauh lebih kecil daripada dugaan &
  Ambang terlalu longgar sehingga orang berbeda tergabung. Naikkan ambang, dan
  periksa daftar \berkas{tinjau-manual.csv}. \\
Jumlah pelanggan master hampir sama dengan jumlah catatan &
  Ambang terlalu ketat, atau normalisasi nama belum dijalankan. Periksa
  keluaran fungsi \perintah{normalkan}. \\
Versi dimensi bertambah setiap pipeline dijalankan &
  Perbandingan atribut memakai \perintah{<>} sehingga perubahan dari atau
  menuju \perintah{NULL} lolos. Gunakan \perintah{IS DISTINCT FROM},
  sebagaimana Modul 3. \\
Jumlah baris fakta berlipat setelah jalan kedua &
  Fakta dimuat tanpa menghapus batch yang sama lebih dahulu. Terapkan
  hapus-lalu-muat per batch. \\
Transaksi terlambat menunjuk harga hari ini &
  Pencarian kunci memakai \perintah{baris\_kini}. Ganti dengan rentang tanggal
  transaksi. \\
\perintah{UnicodeDecodeError} saat membaca sumber &
  Encoding sumber berbeda. Nyatakan encoding secara eksplisit saat koneksi,
  jangan bergantung pada nilai asali sistem. \\
\bottomrule
\end{tabular}
\end{table}

\section{Ringkasan}

Modul ini memperlihatkan bahwa ETL bukan pemindahan data, melainkan
\textbf{rekonsiliasi makna}. Tiga sumber yang sama-sama benar menurut
sistemnya sendiri tidak otomatis dapat digabungkan, dan pekerjaan menyatukannya
sebagian besar adalah keputusan --- satuan mana yang dipakai, kode mana yang
menang, dan dua catatan ini orang yang sama atau bukan.

Tiga gagasan perlu dibawa ke Modul 8. \emph{Pertama}, batas antartahap harus
dijaga: staging menampung apa adanya, dan pembersihan yang dilakukan terlalu
awal menghapus bukti yang diperlukan untuk rekonsiliasi. \emph{Kedua},
keputusan identitas dapat salah, dan kesalahan yang menggabungkan dua orang
jauh lebih berbahaya daripada yang memisahkan satu orang --- karena yang
pertama nyaris mustahil ditemukan kembali. \emph{Ketiga}, pipeline yang tidak
idempoten membuat setiap kegagalan menuntut pembersihan manual, dan kegagalan
adalah hal biasa.

Aturan transformasi yang Anda tulis sebagai \perintah{dict} hari ini dan tabel
\perintah{audit\_muat} yang Anda isi menjadi bahan langsung Modul 8: keduanya
adalah \emph{metadata}, dan Modul 8 mengubah transformasi menjadi kode yang
mendokumentasikan dirinya sendiri beserta graf lineage-nya.
